\scanpage{79}
\diagram{\begin{tikzcd}[ampersand replacement=\&]X\arrow[d]\\Y\end{tikzcd}}
\figure{\begin{tikzpicture}[scale=.78]
\draw (0,1.58)..controls (-1.78,1.49) and (-1.65,-1.55)..(.05,-1.55)..controls (2.12,-1.51) and (1.95,1.55)..(0,1.58);
\draw (-1.44,-.10)..controls (-.4,-.21) and (.58,-.15)..(1.51,-.23);
\draw (.06,1.59)..controls (-.13,.92) and (.08,.25)..(.14,-.15)--(.16,-1.42);
\draw (.08,1.43)..controls (.81,1.32) and (.50,.14)..(.14,-.15);
\fill (.14,-.15) circle (.09);
\node[above] at (.16,1.64){$V(\pi)$};\node[right] at (1.52,-.24){$V(f)$};
\node[left] at (-1.48,1.08){$V\,\bigl\{$};
\draw (-1.44,-4.94)--(1.39,-5.02); \fill (-.13,-4.97) circle (.08);
\end{tikzpicture}}

Soit $\pi\in\unsure{\mathfrak r|r|\mathfrak m}\Lambda$, $\pi\ne0$ donc $\pi$
non div. de 0, \deleted{Alors \illegible{1}} \inserted{Posons}
\[
\deleted{H^{i}_{\unsure{*}}\quad\unsure{V'|V}=X_f}
\]
\[
V=X_f,\quad U=X_{f\pi}=V_{\pi},
\]
\deleted{Alors}
\[
Z=V-U=V\cap V(\pi)=V(\pi)_f,\quad \text{Alors}
\]
\diagram{\begin{tikzcd}[ampersand replacement=\&,column sep=large]
H_Z^i(V,F)\arrow[r]\& H^i(V,F)\arrow[r]\& H^i(U,F)\arrow[ll,bend right=24]
\end{tikzcd}}
et il suffit de prouver que les extrêmes
sont nuls pour $i\geq n+1$. Or $U$ est affine de dim
$\leq n-1$, \inserted{ponct. de t.f.} sur le corps des fractions $K$ de $\Lambda$,
\margin{\#} donc (\unsure{réf|ilf}) $K$ est de dim coh $\leq1$, et
par suite $U$ est de $\ell$-dim ch. $\leq(n-1)+1=n$,
\deleted{D'autre part} ce qui dispose du dernier terme.
Pour le premier, on a
\[
H_Z^*(V,F)\Longleftarrow H^p(Z,\overbrace{\underline H_Z^q(F)})
\]
Or \deleted{je dis que} comme $Z=V(\pi)_f$, \inserted{est sur $k(\Lambda)$,} on
dispose de la dim ch. de $Z$ (qui est
$\leq n-1$) et de ses sous-schémas fermés.

Les termes $E_2^{pq}$ pour $q\ne0,1$ ne donnent
pas de contribution aux $H_{\unsure{?|Z}}^i(V,F)$ pour $i\geq n+1$.
\deleted{Reste} Or je dis que de façon générale
% archive sheet number: 78
